\documentclass[10pt]{article}
\usepackage[margin=0.9in]{geometry}
\usepackage{amsmath,amssymb,booktabs,microtype,xcolor,fancyhdr}
\usepackage[hidelinks]{hyperref}
\definecolor{navy}{RGB}{24,55,95}
\pagestyle{fancy}
\fancyhf{}
\lhead{$\operatorname{Spin}(7)$ on $S_8$}
\rhead{Proof and exact verification}
\cfoot{\thepage}
\setlength{\headheight}{14pt}
\newcommand{\Sg}{\mathcal S}
\newcommand{\CT}{\operatorname{CT}}
\title{\color{navy}\textbf{The compact $\operatorname{Spin}(7)$ sigma-MGF on $S_8$}\\[1mm]
\large Proof and exact degree-four matrix verification}
\author{Computational proof companion}
\date{17 July 2026}

\begin{document}
\maketitle

\begin{abstract}
For the complexification of the real eight-dimensional spin representation,
we prove the compact Molien--Weyl formula, construct explicit maximal-torus
matrices, and verify all coefficients through total degree four by exact
constant terms.  In the Schur basis the answer is
$1+h_2+h_4+s_{(2,2)}+e_4+O_{\ge6}$.  The coefficient of
$m_{(1,1,1,1)}$ is therefore $4$, with the extra invariant supplied by the
Cayley alternating four-form.
\end{abstract}

\section{The coefficient formula}
Let $K$ be compact, let $V$ be a finite-dimensional unitary $K$-module, and
write $h_d(\rho(g))$ for the complete homogeneous character of degree $d$.
Expanding one determinant at a time gives the formal identity
\[
 \prod_{a\ge1}\det(I-t_a\rho(g))^{-1}
 =\prod_{a\ge1}\sum_{d\ge0}h_d(\rho(g))t_a^d.
\]
The coefficient of the monomial symmetric function $m_\tau$, for
$\tau=(\tau_1,\ldots,\tau_r)$, is consequently
\[
 \int_K\prod_{j=1}^r h_{\tau_j}(\rho(g))\,dg
 =\dim\left(\bigotimes_{j=1}^r\operatorname{Sym}^{\tau_j}V\right)^K.
\]
This proves the proposed formula
\[
 \boxed{\Sg_{K,V}=\int_K\prod_a\det(I-t_a\rho(g))^{-1}\,dg
 =\sum_\tau
 \dim\left(\bigotimes_j\operatorname{Sym}^{\tau_j}V\right)^K m_\tau.}
\]
Equivalently, the Cauchy identity gives the Schur-basis expansion
\[
 \Sg_{K,V}=\sum_\lambda\dim(\mathbb S_\lambda V)^K s_\lambda.
\]
The notebook checks both bases independently.

\section{Matrix model and Weyl constant term}
Use type $B_3$ with orthonormal coordinates $e_1,e_2,e_3$.  The weights of
$S_8$ are
\[
 \operatorname{Wt}(S_8)=
 \left\{\tfrac12(\epsilon_1e_1+\epsilon_2e_2+\epsilon_3e_3):
 \epsilon_i\in\{\pm1\}\right\}.
\]
For $\theta\in\mathbb R^3$ this constructs the explicit unitary matrix
\[
 M(\theta)=\operatorname{diag}_{\epsilon\in\{\pm1\}^3}
 \exp\!\left(\frac{i}{2}\sum_{j=1}^3\epsilon_j\theta_j\right)\in U(8).
\]
These matrices form the image of a maximal torus of $\operatorname{Spin}(7)$.
Together with the $18$ roots and $|W(B_3)|=48$, the Weyl integration formula
turns the group integral into the exact constant term
\[
 \boxed{\Sg_{\operatorname{Spin}(7),8}
 =\frac1{48}\CT_z\left[
 \prod_{\alpha>0}(1-z^\alpha)(1-z^{-\alpha})
 \prod_{a,\mu\in\operatorname{Wt}(S_8)}(1-t_a z^\mu)^{-1}
 \right].}
\]
Thus it is unnecessary to approximate Haar measure by random matrices.

\section{Predicted low-degree tensors}
The central element $-1\in\operatorname{Spin}(7)$ acts by $-I$ on $S_8$,
so all odd-total-degree coefficients vanish.  The tensor square is
\[
 S_8\otimes S_8=\mathbf1\oplus\mathbf7\oplus\mathbf{21}\oplus\mathbf{35},
 \qquad
 \operatorname{Sym}^2S_8=\mathbf1\oplus\mathbf{35},\quad
 \bigwedge^2S_8=\mathbf7\oplus\mathbf{21}.
\]
It is multiplicity-free, hence
\[
 \dim(S_8^{\otimes4})^{\operatorname{Spin}(7)}
 =\dim\operatorname{End}_{\operatorname{Spin}(7)}(S_8\otimes S_8)=4.
\]
Three tensors are the metric pairings.  The fourth is the Cayley form in
$\bigwedge^4S_8^*$.  In Schur notation this yields
\[
 \boxed{\Sg_{\operatorname{Spin}(7),8}
 =1+h_2+h_4+s_{(2,2)}+e_4+O_{\ge6}.}
\]
Changing to the monomial basis gives, through degree four,
\[
\begin{array}{c|rrrrrrr}
\tau&(2)&(1,1)&(4)&(3,1)&(2,2)&(2,1,1)&(1,1,1,1)\\ \hline
[m_\tau]&1&1&1&1&2&2&4
\end{array}
\]
Every coefficient of total degree one or three is zero.

\section{Independent exact verification}
The marimo notebook generates the $B_3$ root system and the eight-weight
orbit from the Cartan matrix.  It then performs two distinct checks.
First, for a concrete $M(\theta)$ it compares
\[
 \prod_a\det(I-t_aM(\theta))^{-1}
 \quad\text{with}\quad
 \prod_{a,\mu}(1-t_a e^{i\langle\mu,\theta\rangle})^{-1}.
\]
Second, it forms the characters $h_{\tau_1}\cdots h_{\tau_r}$ directly from
the weights and applies the one-sided Weyl-denominator identity
\[
 \prod_{\alpha>0}(1-z^{-\alpha})
 =\sum_{w\in W}\det(w)z^{w\rho-\rho}.
\]
The resulting constant terms are exact integers and do not use the quoted
tensor-square decomposition as input.

The exact output is
\[
\begin{array}{c|rrrrrrr}
\tau&(2)&(1,1)&(4)&(3,1)&(2,2)&(2,1,1)&(1,1,1,1)\\ \hline
\text{proposed}&1&1&1&1&2&2&4\\
\text{Weyl check}&1&1&1&1&2&2&4
\end{array}
\]
with every omitted partition of degree at most four returning zero.  In the
Schur basis the nonzero multiplicities are
\[
 [s_{(2)}]=[s_{(4)}]=[s_{(2,2)}]=[s_{(1,1,1,1)}]=1.
\]

\paragraph{Reproducibility.}
Run the marimo notebook \texttt{marimo\_notebooks/lie\_group\_spin7.py}.  Its shared
exact-arithmetic core is in the parent folder.  The test suite also executes
this degree-four case noninteractively.

\end{document}
